\documentclass[12pt]{article}
\usepackage[utf8]{inputenc}
\usepackage[T1]{fontenc}
\usepackage{hyperref}
\title{Waxianghua Language Resources for K-12 Curriculum Development}
\author{Compiled by Claude Sonnet 5.5 (claude-sonnet-5-5) for Andr\'as Kornai}
\date{}
\begin{document}
\maketitle

\section{Introduction}
Waxianghua, also called Waxiang or Xianghua (ISO 639-3 wxa; Glottocode waxi1236), is spoken in western Hunan, China. It is mainly spoken in the Wuling Mountains around Yuanling, Chunxi, Jishou, Guzhang and Dayong. Estimates run from about 300,000 to about 400,000 speakers. Ethnologue and Glottolog place it in Sino-Tibetan (Sinitic). Its classification is disputed. Wikipedia treats it as a divergent Chinese variety. Sagart argues that it split early from Old Chinese. Others have proposed a link to Miao.

\section{Primary Educational Resources}
No primers, textbooks or school programmes for Waxianghua appeared in the searches.

\section{Linguistic and Reference Materials}
\begin{itemize}
\item Wikipedia article on Waxiang Chinese. \url{https://en.wikipedia.org/wiki/Waxiang_Chinese}
\item Wikipedia article on the Waxiang people. \url{https://en.wikipedia.org/wiki/Waxiang_people}
\item Wikipedia article on the Macro-Bai languages, which discusses the classification debate. \url{https://en.wikipedia.org/wiki/Macro-Bai_languages}
\item Wu Yunji and Shen Ruiqing (2010), a report on the Guzhang dialect. It is mentioned in search results, but I saw no URL, so it is not in the JSON.
\item Numerals page for Waxianghua from the MPI-EVA Chan numerals database. \url{https://lingweb.eva.mpg.de/channumerals/Waxianghua.htm}
\item Word-list page for Waxianghua (Guzhang) in the Austronesian Basic Vocabulary Database. \url{https://abvd.eva.mpg.de/austronesian/language.php?id=973}
\end{itemize}

\section{Religious and Community Resources}
Joshua Project has a profile for the Waxiang Chinese in China. \url{https://joshuaproject.net/people_groups/18483} Asia Harvest has a profile of the Han Chinese Waxiang. \url{https://people-groups.asiaharvest.org/people-group-profiles/hunan/han-chinese-waxiang.pdf} No scripture translation or community organisation was found.

\section{Digital and Audio Resources}
Verbix hosts a page on Waxianghua. \url{https://docs.verbix.com/Languages/Waxianghua} I did not verify what it contains. No audio, video or software was found in the searches.

\section{Curriculum Development Context}
Little published material exists. The sources seen do not show whether Waxianghua is taught in schools or whether there is a standard writing system. A curriculum would have to start from documentation, and it would need local speakers, Hunan educators and linguists. The unsettled classification and the dominance of Mandarin in education are the main constraints.
\end{document}